
\SGARefTarget{sga3:xi:section:5}\section*{5. First Corollaries of the
Representability Theorem}
\addcontentsline{toc}{section}{5. First Corollaries of the
Representability Theorem}

Let \(H\) and \(G\) be as in \SGARefLink{sga3:xi:corollary:4:2}{4.2}.  Put
\SGARefTarget{sga3:xi:display:5:hom-scheme}
\[
  M=\operatorname{Hom}_{S\text{-}\mathrm{gr}}(H,G),
\]
which, by \SGARefLink{sga3:xi:corollary:4:2}{4.2}, is a smooth \(S\)-prescheme separated over \(S\).
Notice that \(G\) acts on the functor
\(M=\operatorname{Hom}_{S\text{-}\mathrm{gr}}(H,G)\) by
\SGARefTarget{sga3:xi:display:5:conjugation-action}
\[
  (g,u)\longmapsto\operatorname{int}(g)\circ u.
\]
Hence there is a canonical morphism
\SGARefTarget{sga3:xi:display:5:x}
\[
  \tag{x}
  G\times_S M\longrightarrow M\times_S M,
\]
whose components are \SGARefLink{sga3:xi:display:5:conjugation-action}{the preceding morphism}
\(G\times_S M\to M\) and the second projection
\(G\times_S M\to M\).

\medskip
\noindent\SGARefTarget{sga3:xi:corollary:5:1}\textbf{Corollary 5.1.} ---
\textit{The preceding morphism \(\mathrm{\SGARefLink{sga3:xi:display:5:x}{(x)}}\) is smooth.}

This follows from \SGARefLink{sga3:xi:corollary:4:2}{4.2} and \SGARefLink{sga3:xi:corollary:2:3}{2.3}.  This statement is equivalent to the
following one.

\medskip
\noindent\SGARefTarget{sga3:xi:corollary:5:2}\textbf{Corollary 5.2.} ---
\textit{Let \(u_1,u_2:H\rightrightarrows G\) be two homomorphisms of
\(S\)-groups.  Then the subfunctor
\(\operatorname{Transp}(u_1,u_2)\) of \(G\) (cf.\ \SGARefLink{sga3:xi:corollary:2:4}{2.4}) is
representable by a closed subprescheme of \(G\), smooth over \(S\).}

It remains only to prove that
\(\operatorname{Transp}(u_1,u_2)\to G\) is indeed a closed
immersion.  This follows because it is the equalizer of a pair of
morphisms \(G\rightrightarrows M\), namely
\[
  g\longmapsto\operatorname{int}(g)u_1
  \qquad\text{and}\qquad
  g\longmapsto u_2,
\]
the latter being the ``constant'' morphism, and because \(M\) is
separated over \(S\).  In particular:

\medskip
\noindent\SGARefTarget{sga3:xi:corollary:5:3}\textbf{Corollary 5.3.} ---
\textit{Let \(u:H\to G\) be a morphism of \(S\)-groups.  Then
\[
  \operatorname{Centr}_G(u)=\operatorname{Transp}(u,u)
\]
is representable by a closed group subprescheme of \(G\), smooth
over \(S\).  Moreover, the quotient of \(G\) by
\(\operatorname{Centr}_G(u)\) is representable by an open
subprescheme of \(M\).}

It remains to prove the last assertion.  The morphism
\SGARefTarget{sga3:xi:display:5:orbit-map}
\[
  g\longmapsto\operatorname{int}(g)\circ u
\]
from \(G\) to \(M\) is smooth and of finite type by \SGARefLink{sga3:xi:corollary:5:2}{5.2}, hence is an
open morphism (EGA IV, 6.6).  Let \(U\) denote its image, equipped
with the structure induced from \(M\).  The induced morphism
\(G\to U\) is smooth, surjective, and of finite type, and is
therefore a covering for the faithfully flat quasi-compact
topology.  Moreover, it is clear that \SGARefLink{sga3:xi:display:5:orbit-map}{the preceding morphism}
\(G\to M\) makes \(G\) a formally principal homogeneous sheaf under
\(\operatorname{Centr}_G(u)_M\).  It follows that the sheaf
\(G/\operatorname{Centr}_G(u)\) is represented by \(U\).

\medskip
\noindent\SGARefTarget{sga3:xi:corollary:5:4}\textbf{Corollary 5.4.} ---
\textit{Let \(u_1,u_2:H\rightrightarrows G\) be two homomorphisms of
\(S\)-groups.  The following conditions are equivalent:}

\begin{enumerate}
\item[(i)] \SGARefTarget{sga3:xi:corollary:5:4:item:i}
\textit{\(u_1\) and \(u_2\) are conjugate locally for the étale
topology.}

\item[(i bis)] \SGARefTarget{sga3:xi:corollary:5:4:item:i-bis}
\textit{\(u_1\) and \(u_2\) are conjugate locally for the faithfully
flat quasi-compact topology.}

\item[(ii)] \SGARefTarget{sga3:xi:corollary:5:4:item:ii}
\textit{For every \(s\in S\), let \(\overline{s}\) denote the
spectrum of an algebraic closure of \(\kappa(s)\).  The morphisms
\[
  u_{1,\overline{s}},u_{2,\overline{s}}:
  H_{\overline{s}}\rightrightarrows G_{\overline{s}}
\]
are conjugate by an element of
\(G\bigl(\kappa(\overline{s})\bigr)\).}

\item[(ii bis)] \SGARefTarget{sga3:xi:corollary:5:4:item:ii-bis}
\textit{The structural morphism
\(\operatorname{Transp}(u_1,u_2)\to S\) is surjective.}

\item[(iii)] \SGARefTarget{sga3:xi:corollary:5:4:item:iii}
\textit{\(\operatorname{Transp}(u_1,u_2)\) is a torsor under the
action of the smooth \(S\)-group prescheme of finite type
\(\operatorname{Centr}_G(u_1)\).}
\end{enumerate}

The implications \SGARefLink{sga3:xi:corollary:5:4:item:i}{(i)} \(\Rightarrow\) \SGARefLink{sga3:xi:corollary:5:4:item:i-bis}{(i bis)} and
\SGARefLink{sga3:xi:corollary:5:4:item:ii}{(ii)} \(\Rightarrow\) \SGARefLink{sga3:xi:corollary:5:4:item:ii-bis}{(ii bis)} are immediate, the latter by the
Nullstellensatz.  The implication \SGARefLink{sga3:xi:corollary:5:4:item:i-bis}{(i bis)} \(\Rightarrow\) \SGARefLink{sga3:xi:corollary:5:4:item:ii}{(ii)}
follows from the ``principle of finite extension'' (EGA IV, 9).
Moreover, \SGARefLink{sga3:xi:corollary:5:4:item:ii-bis}{(ii bis)} \(\Rightarrow\) \SGARefLink{sga3:xi:corollary:5:4:item:iii}{(iii)}: since
\(\operatorname{Transp}(u_1,u_2)\) is smooth over \(S\), it is flat
over \(S\), and since it is of finite type, it is quasi-compact over
\(S\).  Thus it is faithfully flat and quasi-compact over \(S\) if
and only if its structural morphism is surjective.  On the other
hand, it is formally principal homogeneous under
\(\operatorname{Centr}_G(u_1)\), which is faithfully flat and
quasi-compact over \(S\).  The last condition is therefore also
equivalent to saying that \(\operatorname{Transp}(u_1,u_2)\) is a
torsor under \(\operatorname{Centr}_G(u_1)\), in the faithfully flat
quasi-compact topology.  Finally, \SGARefLink{sga3:xi:corollary:5:4:item:iii}{(iii)} \(\Rightarrow\) \SGARefLink{sga3:xi:corollary:5:4:item:i}{(i)} follows
from the smoothness of \(\operatorname{Transp}(u_1,u_2)\) over \(S\)
and from Hensel's lemma in the form \SGARefLink{sga3:xi:proposition:1:10}{1.10}.

\medskip
\noindent\SGARefTarget{sga3:xi:remark:5:5}\textbf{Remark 5.5.} ---
Fix \(u_1=u:H\to G\).  Consider the functor
\SGARefTarget{sga3:xi:display:5:5:functor}
\[
  (\mathbf{Sch})^\circ/S\longrightarrow(\mathbf{Ens})
\]
that assigns to every \(T\) over \(S\) the set of homomorphisms of
\(T\)-groups \(u_2:H_T\to G_T\) that are locally conjugate, for the
étale topology, to \(u_T:H_T\to G_T\).  \SGARefLink{sga3:xi:display:5:5:functor}{This functor} is represented
precisely by the open subprescheme of \(M\), isomorphic to
\(G/\operatorname{Centr}_G(u)\), considered in \SGARefLink{sga3:xi:corollary:5:3}{5.3}.

Let us sketch the variants of the preceding results obtained by
applying \SGARefLink{sga3:xi:theorem:4:1}{4.1} instead of \SGARefLink{sga3:xi:corollary:4:2}{4.2}.  Thus let \(G\) be a group prescheme
smooth and affine over \(S\), and now let \(M\) denote the smooth
\(S\)-prescheme separated over \(S\) that represents the functor
considered in \SGARefLink{sga3:xi:theorem:4:1}{4.1}.  Again \(G\) acts on \(M\):
\SGARefTarget{sga3:xi:display:5:subgroup-conjugation-action}
\[
  (g,H)\longmapsto\operatorname{int}(g)(H).
\]
Hence, as above, there is a morphism
\SGARefTarget{sga3:xi:display:5:x-bis}
\[
  \tag{x bis}
  G\times_S M\longrightarrow M\times_S M.
\]

\medskip
\noindent\SGARefTarget{sga3:xi:corollary:5:1:bis}\textbf{Corollary 5.1 bis.} ---
\textit{\SGARefLink{sga3:xi:display:5:x-bis}{The preceding morphism} is smooth.}

This follows from \SGARefLink{sga3:xi:theorem:4:1}{4.1} and \SGARefLink{sga3:xi:corollary:2:3:bis}{2.3 bis}.  We again conclude:

\medskip
\noindent\SGARefTarget{sga3:xi:corollary:5:2:bis}\textbf{Corollary 5.2 bis.} ---
\textit{Let \(H_1,H_2\) be two multiplicative-type subgroups of
\(G\).  Then the subfunctor
\(\operatorname{Transp}_G(H_1,H_2)\) of \(G\) is representable by a
closed subprescheme of \(G\), smooth over \(S\).}

In particular:

\medskip
\noindent\SGARefTarget{sga3:xi:corollary:5:3:bis}\textbf{Corollary 5.3 bis.} ---
\textit{Let \(H\) be a multiplicative-type subgroup of \(G\).  Then
the group subfunctor \(\operatorname{Norm}_G(H)\) of \(G\) is
representable by a closed subprescheme of \(G\), smooth over \(S\).
Moreover, the quotient \(G/\operatorname{Norm}_G(H)\) is
representable by an open subprescheme of \(M\).}

\medskip
\noindent\SGARefTarget{sga3:xi:corollary:5:4:bis}\textbf{Corollary 5.4 bis.} ---
\textit{Let \(H_1,H_2\) be two multiplicative-type subgroups of
\(G\).  The following conditions are equivalent:}

\begin{enumerate}
\item[(i)] \SGARefTarget{sga3:xi:corollary:5:4:bis:item:i}
\textit{\(H_1\) and \(H_2\) are conjugate locally for the étale
topology.}

\item[(i bis)] \SGARefTarget{sga3:xi:corollary:5:4:bis:item:i-bis}
\textit{\(H_1\) and \(H_2\) are conjugate locally for the faithfully
flat quasi-compact topology.}

\item[(ii)] \SGARefTarget{sga3:xi:corollary:5:4:bis:item:ii}
\textit{For every \(s\in S\), let \(\overline{s}\) denote the
spectrum of an algebraic closure of \(\kappa(s)\).  The subgroups
\(H_{1,\overline{s}}\) and \(H_{2,\overline{s}}\) of
\(G_{\overline{s}}\) are conjugate by an element of
\(G\bigl(\kappa(\overline{s})\bigr)\).}

\item[(ii bis)] \SGARefTarget{sga3:xi:corollary:5:4:bis:item:ii-bis}
\textit{The structural morphism
\(\operatorname{Transp}(H_1,H_2)\to S\) is surjective.}

\item[(iii)] \SGARefTarget{sga3:xi:corollary:5:4:bis:item:iii}
\textit{\(\operatorname{Transp}(H_1,H_2)\) is a principal
homogeneous bundle under the action of the smooth \(S\)-group
prescheme of finite type \(\operatorname{Norm}_G(H_1)\).}
\end{enumerate}

\medskip
\noindent\SGARefTarget{sga3:xi:remark:5:5:bis}\textbf{Remark 5.5 bis.} ---
\SGARefLink{sga3:xi:remark:5:5}{Remark 5.5} carries over to the present case as well.

\medskip
\noindent\SGARefTarget{sga3:xi:remark:5:6}\textbf{Remark 5.6.} ---
To prove \SGARefLink{sga3:xi:corollary:5:2}{5.2}, and hence also the first assertion of \SGARefLink{sga3:xi:corollary:5:3}{5.3} and the
assertions of \SGARefLink{sga3:xi:corollary:5:4}{5.4}, the reference to \SGARefLink{sga3:xi:corollary:4:2}{4.2} may be replaced by a
reference to VIII, \SGARefLink{sga3:viii:theorem:6:4}{6.4}, whose proof is much easier.  This also shows
that the hypothesis that \(G\) be affine over \(S\) is unnecessary
there.  Moreover, in \SGARefLink{sga3:xi:section:6}{Section 6} we shall show how a variant of this
method extends these results to certain groups \(H\) more general
than groups of multiplicative type.  The same observations apply to
the variants \SGARefLink{sga3:xi:corollary:5:2:bis}{5.2 bis} and following.  By contrast, the result \SGARefLink{sga3:xi:theorem:5:8}{5.8}
below makes essential use of all the hypotheses imposed---in
particular, that \(G\) be affine and smooth over \(S\), and that
\(H\) be of multiplicative type---and the lecturer knows no proof
other than one using the representability theorems \SGARefLink{sga3:xi:theorem:4:1}{4.1} or
\SGARefLink{sga3:xi:corollary:4:2}{4.2}.\footnote{The situation has changed since this text was written;
cf.\ Exposés XV and XIX, no.~6.}

\medskip
\noindent\SGARefTarget{sga3:xi:paragraph:5:7}\textbf{5.7.}
Since the morphism \(\mathrm{\SGARefLink{sga3:xi:display:5:x}{(x)}}\), respectively
\SGARefLink{sga3:xi:display:5:x-bis}{\(\mathrm{(x\ bis)}\)}, is smooth and hence open, its image is open.
Let \(R\) be this image, equipped with the structure induced from
\(M\times_S M\).  One checks readily that \(R\) is an equivalence
relation on \(M\), namely the relation made explicit in \SGARefLink{sga3:xi:corollary:5:4}{5.4},
respectively \SGARefLink{sga3:xi:corollary:5:4:bis}{5.4 bis}.  It would be interesting to know whether the
quotient sheaf \(M/R\), which is formally étale over \(S\), is
representable; if so, it is represented by a prescheme étale over
\(S\).  This is the case, for example, when \(S\) is the spectrum of
a field.  Notice also that \(R\) need not be closed in
\(M\times_S M\), even when \(G\) is quasi-finite over \(S\).  Thus,
when \(M/R\) is representable, it may be nonseparated over \(S\).

\medskip
\noindent\SGARefTarget{sga3:xi:theorem:5:8}\textbf{Theorem 5.8.} ---
\textit{Let \(S\) be a prescheme, let \(G\) be a smooth affine
\(S\)-prescheme, and let \(s\in S\).  Then:}

\begin{enumerate}
\item[(a)] \SGARefTarget{sga3:xi:theorem:5:8:item:a}
\textit{For every multiplicative-type subgroup \(H_0\) of \(G_s\),
there exist an étale morphism \(S'\to S\), a point \(s'\in S'\)
lying over \(s\) such that the residue-field extension
\(\kappa(s')/\kappa(s)\) is trivial, and a multiplicative-type
subgroup \(H'\) of \(G'=G\times_S S'\), such that}
\SGARefTarget{sga3:xi:display:5:8:item:a:fiber}
\[
  H'_{s'}=H_0\otimes_{\kappa(s)}\kappa(s').
\]

\item[(b)] \SGARefTarget{sga3:xi:theorem:5:8:item:b}
\textit{For every group homomorphism \(u_0:H_s\to G_s\), where
\(H\) is a multiplicative-type \(S\)-group prescheme of finite
type, there exist an étale morphism \(S'\to S\), a point \(s'\in S'\)
lying over \(s\) such that \(\kappa(s')/\kappa(s)\) is trivial, and
a group homomorphism
\[
  u':H'\longrightarrow G',
  \qquad
  H'=H\times_S S',
  \quad
  G'=G\times_S S',
\]
such that}
\SGARefTarget{sga3:xi:display:5:8:item:b:fiber}
\[
  u'_{s'}:H'_{s'}\longrightarrow G'_{s'}
  \quad\text{equals}\quad
  u_0\otimes_{\kappa(s)}\kappa(s').
\]
\end{enumerate}

This follows from \SGARefLink{sga3:xi:theorem:4:1}{4.1}, respectively \SGARefLink{sga3:xi:corollary:4:2}{4.2}, and from Hensel's lemma in
the form \SGARefLink{sga3:xi:proposition:1:10}{1.10}.

\medskip
\noindent\SGARefTarget{sga3:xi:proposition:5:9}\textbf{Proposition 5.9.} ---
\textit{Let \(S\) be a prescheme, let \(G\) be a smooth affine
\(S\)-prescheme, and let \(H\) be a multiplicative-type subgroup of
\(G\).  Consider}
\[
  \operatorname{Centr}_G(H)\lhook\joinrel\longrightarrow
  \operatorname{Norm}_G(H),
\]
\textit{which, by \SGARefLink{sga3:xi:corollary:5:3}{5.3} and \SGARefLink{sga3:xi:corollary:5:3:bis}{5.3 bis}, are closed group subpreschemes of
\(G\), smooth over \(S\).  The first is then an open and closed
subprescheme of the second, and the quotient sheaf}
\SGARefTarget{sga3:xi:display:5:9:weyl-quotient}
\[
  W_G(H)=
  \operatorname{Norm}_G(H)\big/\operatorname{Centr}_G(H)
\]
\textit{is represented by an open group subprescheme of
\(\operatorname{Aut}_{S\text{-}\mathrm{gr}}(H)\).  It is therefore
a quasi-finite, étale, and separated \(S\)-group prescheme.}

Indeed, consider the evident homomorphism
\SGARefTarget{sga3:xi:display:5:9:theta}
\[
  \theta:\operatorname{Norm}_G(H)\longrightarrow
  \operatorname{Aut}_{S\text{-}\mathrm{gr}}(H),
\]
whose kernel is, by definition, \(\operatorname{Centr}_G(H)\).
Since \(\operatorname{Aut}_{S\text{-}\mathrm{gr}}(H)\) is
represented by an étale \(S\)-group prescheme separated over \(S\)
(X, \SGARefLink{sga3:x:corollary:5-10}{5.10}), its unit section is an open and closed immersion.  Its
inverse image under \(\theta\) is therefore an open and closed
subgroup of \(\operatorname{Norm}_G(H)\).

Moreover, \(\theta\) is smooth.  This follows formally from the
definitions and from the facts that \(\operatorname{Norm}_G(H)\) is
smooth over \(S\) and that
\(\operatorname{Aut}_{S\text{-}\mathrm{gr}}(H)\) is étale over
\(S\).  As in \SGARefLink{sga3:xi:corollary:5:3}{5.3}, the image of \(\theta\) is consequently an open
subprescheme \(U\) of
\(\operatorname{Aut}_{S\text{-}\mathrm{gr}}(H)\), and, with the
induced structure, \(U\) represents the quotient sheaf
\(\operatorname{Norm}_G(H)/\operatorname{Centr}_G(H)\).  The latter
is therefore étale and separated over \(S\).  It is also quasi-finite
over \(S\), since it is quasi-compact over \(S\), being the image of
\(\operatorname{Norm}_G(H)\), which is quasi-compact.  This proves
\SGARefLink{sga3:xi:proposition:5:9}{5.9}.

\medskip
\noindent\SGARefTarget{sga3:xi:corollary:5:10}\textbf{Corollary 5.10.} ---
\textit{For every \(s\in S\), put}
\SGARefTarget{sga3:xi:display:5:10:w}
\[
  w(s)=
  \operatorname{rank}
  \bigl(
    \operatorname{Norm}_{G_s}(H_s)/
    \operatorname{Centr}_{G_s}(H_s)
  \bigr).
\]
\textit{This is also the index of
\(\operatorname{Centr}_{G(\overline{k})}(H(\overline{k}))\) in
\(\operatorname{Norm}_{G(\overline{k})}(H(\overline{k}))\), where
\(\overline{k}\) is an algebraic closure of \(\kappa(s)\).  Then the
function \(s\mapsto w(s)\) is lower semicontinuous.  It is constant
in a neighborhood of \(s\) if and only if \(W_G(H)\) is finite over
\(S\) in a neighborhood of \(s\).}

Indeed, for every \(S\)-prescheme \(W\) that is étale, of finite
type, and separated over \(S\), the function
\[
  s\longmapsto
  \operatorname{rank}[W_s:\kappa(s)]
\]
is lower semicontinuous, and it is constant in a neighborhood of
\(s\) if and only if \(W\) is finite over \(S\) in a neighborhood of
\(s\).  This fact was noted in SGA 1, I, 10.9; its proof, which
presents no difficulty, is given in EGA IV.\footnote{EGA IV, 15.5.1
and 18.10.7.}

\medskip
\noindent\SGARefTarget{sga3:xi:remark:5:11}\textbf{Remark 5.11.}%
\footnote{Editor's note: For a generalization to the nonaffine case,
see M.~Raynaud, \textit{Faisceaux amples sur les schémas en groupes
et les espaces homogènes}, Lecture Notes in Mathematics 119 (1970),
IX.2.8.} ---
Let \(G\) be a group prescheme affine and smooth over \(S\), and let
\(H\) be a multiplicative-type subgroup.  Then \SGARefLink{sga3:xi:corollary:5:3}{5.3} and \SGARefLink{sga3:xi:corollary:5:3:bis}{5.3 bis} imply
that the quotients
\SGARefTarget{sga3:xi:display:5:11:quotients}
\[
  G/\operatorname{Centr}_G(H)
  \qquad\text{and}\qquad
  G/\operatorname{Norm}_G(H)
\]
are preschemes smooth and quasi-affine over \(S\).  Indeed, in both
cases the moduli prescheme \(M\) in which the quotient is embedded
has the property that every open \(U\subset M\) quasi-compact over
\(S\) is quasi-affine over \(S\), by \SGARefLink{sga3:xi:corollary:3:11}{3.11}.

We shall moreover see in the next Exposé\footnote{Editor's note:
Cf.\ Remark XII.5.7.  We mention here the following generalization.
Without assuming \(G\) affine, if \(H\) is a torus and \(S\) is
locally noetherian, \(U\) is affine and smooth by a theorem of
Raynaud (\textit{loc.\ cit.}, IX.2.6 and IX.2.8).  Indeed,
\(G/\operatorname{Norm}_G(H)\) is an open subprescheme of \(M\),
and the largest representable open subfunctor of \(M\) is a disjoint
sum of open subpreschemes smooth and affine over \(S\).} that for
every such \(U\), the schematic closure \(\overline{U}\) of \(U\) in
\(M\) is even affine over \(S\).  Thus the quotients under
consideration are affine over \(S\), provided that the open
subprescheme \(U\) of the corresponding moduli prescheme---of group
homomorphisms from \(H\) to \(G\), respectively of
multiplicative-type subgroups of \(G\)---is closed in \(M\).  This
is the case, for example, when \(H\) is a ``maximal torus'' in the
second case, or when \(S\) is the spectrum of a field; cf.\
XII.5.4.\footnote{Editor's note: This assertion assumes that \(G\)
has locally constant reductive rank.}  I do not know whether the
quotients under consideration are affine over \(S\) in general.

